% ECONOMICS_PROBLEM: {"id": "C26-181", "title": "Weak and many instruments", "jel_code": "C26", "source_id": "OP-0103"}
% !TeX program = xelatex
\documentclass[11pt,a4paper]{article}
\usepackage[margin=23mm]{geometry}
\usepackage{fontspec}
\setmainfont{Latin Modern Roman}
\usepackage{amsmath,amssymb,mathtools}
\usepackage{xcolor,enumitem,booktabs,longtable,array,xurl,hyperref}
\definecolor{muted}{HTML}{53636C}
\hypersetup{colorlinks=true,urlcolor=blue,linkcolor=blue}
\setlength{\parindent}{0pt}
\setlength{\parskip}{5pt}
\newcommand{\R}{\mathbb R}
\newcommand{\E}{\mathbb E}
\newcommand{\N}{\mathbb N}
\newcommand{\ind}{\mathbf 1}
\newcommand{\Var}{\operatorname{Var}}
\newcommand{\Cov}{\operatorname{Cov}}
\newcommand{\supp}{\operatorname{supp}}
\DeclareMathOperator*{\argmin}{arg\,min}
\DeclareMathOperator*{\argmax}{arg\,max}
\newcommand{\entrymeta}[3]{{\sffamily\small\color{muted}\textbf{\texttt{#1}}\quad|\quad #2\par #3\par}}
\newcommand{\Problemref}[1]{\texttt{#1}}
\newcommand{\JELref}[2]{\texttt{#2} (JEL \texttt{#1})}
\begin{document}
\section*{Weak and many instruments}
\textbf{Problem ID:} \texttt{C26-181}\quad \textbf{JEL:} \texttt{C26}\par
% BEGIN ECONOMICS STATEMENT
\label{problem:OP-0103}
\label{atlas:prob:E3}

\noindent\textbf{Original identifier:} E3 (atlas item 103). \textbf{Field:} Econometrics and causal inference. \textbf{Classification:} editorial research family with established restricted solutions; current open status not certified.

\subsection*{Definitions and assumptions}

Consider independent observations $(Y_i,D_i,Z_i)$ for $i=1,\ldots,n$, after any stated exogenous controls have been partialled out. Here $Y_i,D_i\in\mathbb R$ and $Z_i\in\mathbb R^{k_n}$ with $k_n$ potentially growing. A linear structural model is $Y_i=D_i\beta+u_i$, $D_i=Z_i^\top\pi_n+v_i$, with finite stated moments. Under instrument validity $\mathbb E[Z_i u_i]=0$ and normalized nonsingular $\mathbb E[Z_iZ_i^\top]$, relevance may weaken, for example $\pi_n=c/\sqrt n$ for fixed dimension and bounded $c$. Neither a nonzero sample first-stage estimate nor the number of instruments establishes identification.

\subsection*{Formal research question}

Find identification-robust confidence sets and an explicit efficiency frontier when $k_n$ grows, first-stage strength can approach zero, and instrument invalidity is bounded rather than excluded. For a supplied norm $\|\cdot\|_*$ and tolerance $\varepsilon_n\ge0$, define
\[
\mathcal B(P,\varepsilon_n)=
\{b\in\mathbb R:\|\mathbb E_P[Z(Y-Db)]\|_*\le\varepsilon_n\}.
\]
For a specified class $\mathcal P_n$ with moment, dependence and dimensional restrictions, construct $C_n$ with $\liminf_n\inf_{P\in\mathcal P_n}\Pr_P\{\mathcal B(P,\varepsilon_n)\subseteq C_n\}\ge1-\alpha$, where $\alpha\in(0,1)$. Determine necessary growth restrictions on $k_n$, appropriate covariance regularization, and whether selection of instruments on the same data can preserve this guarantee. Compare expected diameter only over submodels where finite informative sets are possible, and establish impossibility results where identification disappears.

\subsection*{Interpretation and scope}

The norm and tolerance must fix instrument scaling; otherwise an invalidity bound changes when an instrument is multiplied by a constant. With zero relevance, $\beta$ need not be point identified and an honest confidence set may be unbounded. Anderson--Rubin-type tests already solve important weak-instrument cases, while many-instrument asymptotics have their own regimes. Arbitrary invalid instruments cannot identify a causal slope without additional restrictions. Clustered sampling changes the experiment and requires enough independent clusters, not simply a large observation count. A causal interpretation further requires a defensible structural equation.

\subsection*{Source audit and status}

[S1] develops local-to-zero asymptotics; [S2] surveys weak identification in IV and GMM; [S3] reviews both theory and practice, including heteroskedastic and dependent settings. Reliable inference with weak instruments is therefore not an entirely unsolved question. The combined growing-dimension, bounded-invalidity and adaptive-selection target is an editorial extension; the cited reviews do not certify its exact 2026 status.

\subsection*{References}

\begin{enumerate}[label={[\textnormal{S}\arabic*]},leftmargin=*]

\item Douglas Staiger and James H. Stock (1997). \emph{Instrumental Variables Regression with Weak Instruments}. Econometrica 65(3), 557–586. \url{https://www.nber.org/papers/t0151}. \textbf{Locator:} Publisher or institutional abstract and bibliographic record. \textbf{Support and limits:} Local-to-zero weak-instrument asymptotics; NBER page verifies the 1997 publication, although the working paper is 1994.

\item James H. Stock, Jonathan H. Wright and Motohiro Yogo (2002). \emph{A Survey of Weak Instruments and Weak Identification in Generalized Method of Moments}. Journal of Business \& Economic Statistics 20(4), 518–529. \url{https://stock.scholars.harvard.edu/publications/survey-weak-instruments-and-weak-identification-generalized-method-moments}. \textbf{Locator:} Publisher or institutional abstract and bibliographic record. \textbf{Support and limits:} Survey of established weak-identification methods; does not establish inference with arbitrary invalid instruments.

\item Isaiah Andrews, James H. Stock and Liyang Sun (2019). \emph{Weak Instruments in Instrumental Variables Regression: Theory and Practice}. Annual Review of Economics 11, 727–753. \url{https://doi.org/10.1146/annurev-economics-080218-025643}. \textbf{Locator:} Publisher or institutional abstract and bibliographic record. \textbf{Support and limits:} Reviews identification-robust inference, including nonhomoskedastic data; the basic question is already partly solved.

\end{enumerate}

\subsection*{Common conventions: Source mathematical conventions}

Unless an entry states otherwise, agent and firm sets are finite. State and action spaces are measurable, strategies and outcome maps are measurable, probability laws are specified, and expectations exist for the targets under discussion. Infinite-horizon discount factors lie in $(0,1)$ and discounted objectives are finite. Norms, tolerances and complexity input sizes must be specified for computational comparisons. Constants or primitives that appear locally are inputs, not empirical facts asserted by this catalogue.

\textbf{Identification.} A statistical model $m\in\mathcal M$ implies an observable law $P_m$ and target $\theta(m)$. At law $P$, its identified set is
\[
\Theta(P;\mathcal M)=\{\theta(m):m\in\mathcal M,\ P_m=P\}.
\]
Point identification holds when this set is a singleton, or equivalently when $P_{m_1}=P_{m_2}$ implies $\theta(m_1)=\theta(m_2)$ for all compatible models. Sharp bounds mean bounds attained, or approached when the boundary is not attained, by compatible models. Writing this set is a definition, not a solution: the substantive task is to establish informative restrictions and derive or estimate the set. An unrestricted-confounding model may give only trivial bounds.

\textbf{Inference.} Identification, estimability and valid uncertainty quantification are separate questions. Fix $\alpha\in(0,1)$. A confidence procedure $C_n$ over a declared law class $\mathcal P$ has asymptotic uniform coverage at level $1-\alpha$ when
\[
\liminf_{n\to\infty}\inf_{P\in\mathcal P}\Pr_P\{\theta(P)\in C_n\}\geq1-\alpha.
\]
For set-identified targets the coverage event must be defined separately, for example coverage of the full identified set. If an entry uses identification-indexed models instead of uniquely indexed laws, the same coverage convention is applied over those models. Pointwise limits do not establish uniform validity, and asymptotic validity does not guarantee finite-sample precision. Sample sizes, dependence structures and weak-identification sequences are part of each application.

\textbf{Counterfactuals.} $Y_i(a)$ is a potential outcome under a specified intervention, including its timing, implementation and versions. It is not $Y_i$ conditional on voluntarily choosing $a$. A full intervention vector permits interference; shorthand scalar treatment notation does not silently impose no interference. Assignment, selection, attrition, anticipation and measurement must be specified. A claim about an instrument's local effect is not automatically a population policy effect.

\textbf{Economic equilibrium.} A counterfactual model includes best responses, constraints, feasibility, market clearing, state transitions and expectations consistent with the stated information. When several equilibria exist, either a declared selector is an input or the target is set-valued. Comparisons across policy regimes must use the stated selection convention. Reduced-form relationships are not presumed invariant after an equilibrium-changing intervention.

\textbf{Optimization and welfare.} Optimization is over the stated feasible set. If attainment is not established, $\sup$ or $\inf$ and $\varepsilon$-optimal choices replace an asserted maximizer or minimizer. For example, with value $V=\sup_{a\in\mathcal A}W(a)<\infty$, an $\varepsilon$-optimal set is $\{a\in\mathcal A:W(a)\geq V-\varepsilon\}$ for $\varepsilon>0$. Benefits, costs, risk and health or environmental outcomes can be aggregated only using declared units and conversion weights. Interpersonal utility weights, the treatment of fiscal transfers and foreign profits, discounting, and the behavioral welfare criterion are explicit inputs. A comparative-static derivative additionally requires existence and appropriate differentiability of the selected counterfactual.

\textbf{Sources.} Local labels [S1], [S2], and so forth restart in every question. Locators identify the relevant abstract, model, section or result at the level actually checked; a bibliographic or abstract-level verification is not represented as inspection of an entire proof. Working papers are distinguished from later publications and changing versions. Source summaries are paraphrased. All URL checks and editorial formulations refer to the audit date stated above.
% END ECONOMICS STATEMENT
\end{document}
